\chapter{Method and Message Reference}
\label{ch:method-reference}

Everything in protocol revision \texttt{2026-07-28}, on a few pages. Where a
chapter covers something properly, it is cross-referenced rather than repeated.

\section{Methods}

\begin{table}[htbp]
\centering
\small
\begin{tabularx}{\textwidth}{@{}lXccl@{}}
\toprule
\thead{Method} & \thead{Purpose} & \thead{Page} & \thead{Cache} & \thead{MRTR} \\
\midrule
\texttt{server/discover}          & Versions, capabilities, identity & no  & yes & no \\
\texttt{tools/list}               & Available tools                  & yes & yes & no \\
\texttt{tools/call}               & Invoke a tool                    & no  & no  & \textbf{yes} \\
\texttt{resources/list}           & Available resources              & yes & yes & no \\
\texttt{resources/templates/list} & Parameterised resource families  & yes & yes & no \\
\texttt{resources/read}           & Resource contents                & no  & yes & \textbf{yes} \\
\texttt{prompts/list}             & Available prompts                & yes & yes & no \\
\texttt{prompts/get}              & Render a prompt                  & no  & no  & \textbf{yes} \\
\texttt{completion/complete}      & Autocomplete an argument         & no  & no  & no \\
\texttt{subscriptions/listen}     & Long-lived notification stream   & no  & no  & no \\
\bottomrule
\end{tabularx}
\caption{The complete core surface. ``Page'' means the operation supports
cursor pagination; ``Cache'' means results carry \texttt{ttlMs} and
\texttt{cacheScope}; ``MRTR'' means the server may answer
\texttt{input\_required}.}
\label{tab:all-methods}
\end{table}

Ten methods. That is the entire core protocol.

\subsection{Tasks extension (\texttt{io.modelcontextprotocol/tasks})}

\begin{table}[htbp]
\centering
\small
\begin{tabularx}{\textwidth}{@{}lX@{}}
\toprule
\texttt{tasks/get}    & Poll a task. Honour \texttt{pollIntervalMs} \\
\texttt{tasks/update} & Supply \texttt{inputResponses}. Must be idempotent \\
\texttt{tasks/cancel} & Request cancellation. Cooperative, not binding \\
\bottomrule
\end{tabularx}
\caption{See Chapter~\ref{ch:tasks}.}
\label{tab:tasks-methods}
\end{table}

\subsection{Removed in this revision}

\begin{table}[htbp]
\centering
\small
\begin{tabularx}{\textwidth}{@{}lX@{}}
\toprule
\thead{Gone} & \thead{Replacement} \\
\midrule
\texttt{initialize} & Per-request \texttt{\_meta} \\
\texttt{notifications/initialized} & nothing \\
\texttt{ping} & nothing; use a transport-level check \\
\texttt{logging/setLevel} & \texttt{io.modelcontextprotocol/logLevel} per request \\
\texttt{resources/subscribe} & \texttt{subscriptions/listen} \\
\texttt{resources/unsubscribe} & close the stream \\
\texttt{notifications/roots/list\_changed} & nothing; Roots is deprecated \\
\texttt{tasks/result}, \texttt{tasks/list} & \texttt{tasks/get} \\
\bottomrule
\end{tabularx}
\caption{Sending any of these to a strictly modern server yields
\texttt{-32601}.}
\label{tab:removed-methods}
\end{table}

\section{Notifications}

\begin{table}[htbp]
\centering
\small
\begin{tabularx}{\textwidth}{@{}lXl@{}}
\toprule
\thead{Notification} & \thead{Meaning} & \thead{Channel} \\
\midrule
\texttt{notifications/progress} & Progress on a request & that request's stream \\
\texttt{notifications/message} & A log line & that request's stream \\
\texttt{notifications/subscriptions/acknowledged} & Subscription accepted & listen stream \\
\texttt{notifications/tools/list\_changed} & Catalogue changed & listen stream \\
\texttt{notifications/prompts/list\_changed} & Prompts changed & listen stream \\
\texttt{notifications/resources/list\_changed} & Resource list changed & listen stream \\
\texttt{notifications/resources/updated} & One resource changed & listen stream \\
\texttt{notifications/cancelled} & Client cancels a request & stdio only \\
\bottomrule
\end{tabularx}
\caption{The channel column is the part people get wrong. Request-scoped
notifications never appear on the listen stream.}
\label{tab:notifications}
\end{table}

\section{The \texttt{\_meta} envelope}

\begin{table}[htbp]
\centering
\small
\begin{tabularx}{\textwidth}{@{}lXcl@{}}
\toprule
\thead{Key} & \thead{Purpose} & \thead{Req.} & \thead{On} \\
\midrule
\texttt{io.modelcontextprotocol/protocolVersion} & Version for this request & \textbf{yes} & requests \\
\texttt{io.modelcontextprotocol/clientCapabilities} & What the client supports & \textbf{yes} & requests \\
\texttt{io.modelcontextprotocol/clientInfo} & Client name and version & no & requests \\
\texttt{io.modelcontextprotocol/logLevel} & Minimum log level & no & requests \\
\texttt{io.modelcontextprotocol/serverInfo} & Server name and version & no & results \\
\texttt{io.modelcontextprotocol/subscriptionId} & Correlates a notification & yes & listen stream \\
\texttt{progressToken} & Opts into progress & no & requests \\
\texttt{traceparent} & W3C trace context & no & requests \\
\texttt{tracestate} & W3C trace state & no & requests \\
\texttt{baggage} & W3C baggage & no & requests \\
\bottomrule
\end{tabularx}
\caption{See Chapter~\ref{ch:architecture}. The last three are an explicit
exception to the reverse-DNS prefix rule.}
\label{tab:meta-reference}
\end{table}

\subsection{Key naming}

Prefixes are reverse DNS with a trailing slash. A prefix whose \emph{second}
label is \texttt{modelcontextprotocol} or \texttt{mcp} is reserved.

\begin{table}[htbp]
\centering
\small
\begin{tabular}{@{}ll@{}}
\toprule
\texttt{io.modelcontextprotocol/x} & reserved \\
\texttt{dev.mcp/x} & reserved \\
\texttt{org.modelcontextprotocol.api/x} & reserved \\
\texttt{com.mcp.tools/x} & reserved \\
\texttt{com.example.mcp/x} & \textbf{yours} \\
\texttt{com.meridian/promptVersion} & \textbf{yours} \\
\bottomrule
\end{tabular}
\caption{The second label is what counts.}
\label{tab:meta-prefixes}
\end{table}

\section{Error codes}

\begin{table}[htbp]
\centering
\small
\begin{tabularx}{\textwidth}{@{}llXl@{}}
\toprule
\thead{Code} & \thead{Name} & \thead{Meaning} & \thead{HTTP} \\
\midrule
\texttt{-32700} & ParseError & Invalid JSON & 400 \\
\texttt{-32600} & InvalidRequest & Bad JSON-RPC framing & 400 \\
\texttt{-32601} & MethodNotFound & Unknown method & 404 \\
\texttt{-32602} & InvalidParams & Bad params, unknown tool, missing resource & 400 \\
\texttt{-32603} & InternalError & Server fault & 500 \\
\texttt{-32020} & HeaderMismatch & Header disagrees with body & 400 \\
\texttt{-32021} & MissingRequiredClientCapability & Client lacks a capability & 400 \\
\texttt{-32022} & UnsupportedProtocolVersion & Version not supported & 400 \\
\bottomrule
\end{tabularx}
\caption{See Chapter~\ref{ch:transport}.}
\label{tab:error-reference}
\end{table}

\begin{table}[htbp]
\centering
\small
\begin{tabularx}{\textwidth}{@{}lX@{}}
\toprule
\thead{Range} & \thead{Rule} \\
\midrule
\texttt{-32000} to \texttt{-32019} & Legacy. Do not allocate; assume no meaning \\
\texttt{-32020} to \texttt{-32099} & \textbf{Reserved for the specification} \\
Outside \texttt{-32768..-32000} & Yours \\
\bottomrule
\end{tabularx}
\caption{Retired and never reissued: \texttt{-32002} (resource not found,
now \texttt{-32602}) and \texttt{-32042} (URL elicitation required).}
\label{tab:error-ranges-ref}
\end{table}

\section{Result types}

\begin{table}[htbp]
\centering
\small
\begin{tabularx}{\textwidth}{@{}llX@{}}
\toprule
\thead{\texttt{resultType}} & \thead{Source} & \thead{Client does} \\
\midrule
\texttt{"complete"} & core & Reads the result \\
\texttt{"input\_required"} & core & Gathers input, retries with a new id \\
\texttt{"task"} & Tasks ext. & Polls \texttt{tasks/get} \\
absent & pre-2026 servers & Treats as \texttt{"complete"} \\
unrecognised & anywhere & \textbf{Treats as invalid} \\
\bottomrule
\end{tabularx}
\caption{See Chapter~\ref{ch:architecture}.}
\label{tab:result-types}
\end{table}

\section{HTTP headers}

\begin{table}[htbp]
\centering
\small
\begin{tabularx}{\textwidth}{@{}lXl@{}}
\toprule
\thead{Header} & \thead{Mirrors} & \thead{Required} \\
\midrule
\texttt{MCP-Protocol-Version} & \texttt{\_meta} protocol version & every POST \\
\texttt{Mcp-Method} & \texttt{method} & every POST \\
\texttt{Mcp-Name} & \texttt{params.name} or \texttt{params.uri} & \texttt{tools/call}, \texttt{prompts/get}, \texttt{resources/read} \\
\texttt{Mcp-Param-\{Name\}} & An \texttt{x-mcp-header} argument & when the schema declares one \\
\texttt{Accept} & & must list both JSON and SSE \\
\texttt{X-Accel-Buffering: no} & & response header on every SSE stream \\
\bottomrule
\end{tabularx}
\caption{Any mismatch between header and body must be rejected with 400 and
\texttt{-32020}. See Chapter~\ref{ch:transport}.}
\label{tab:header-reference}
\end{table}

Values outside printable ASCII, with surrounding whitespace, or resembling the
sentinel, are wrapped: \texttt{=?base64?\{payload\}?=}.

\section{Capabilities}

\begin{table}[htbp]
\centering
\small
\begin{tabularx}{\textwidth}{@{}llX@{}}
\toprule
\thead{Side} & \thead{Capability} & \thead{Settings} \\
\midrule
Server & \texttt{tools} & \texttt{listChanged} \\
Server & \texttt{resources} & \texttt{listChanged}, \texttt{subscribe} \\
Server & \texttt{prompts} & \texttt{listChanged} \\
Server & \texttt{completions} & \\
Both & \texttt{extensions} & map of identifier to settings \\
Client & \texttt{elicitation} & \texttt{form}, \texttt{url} \\
Client & \texttt{sampling} & \emph{deprecated} \\
Client & \texttt{roots} & \emph{deprecated} \\
\bottomrule
\end{tabularx}
\caption{An empty \texttt{elicitation} object means form mode only, for
backwards compatibility.}
\label{tab:capabilities-reference}
\end{table}

\section{The deprecated features registry}

\begin{table}[htbp]
\centering
\small
\begin{tabularx}{\textwidth}{@{}lXl@{}}
\toprule
\thead{Feature} & \thead{Migrate to} & \thead{Earliest removal} \\
\midrule
Roots & Tool parameters, resource URIs, or server config & 2027-07-28 \\
Sampling & Direct provider APIs & 2027-07-28 \\
Logging & \texttt{stderr} on stdio; OpenTelemetry elsewhere & 2027-07-28 \\
Dynamic Client Registration & Client ID Metadata Documents & 2027-07-28 \\
\texttt{includeContext: "thisServer"} / \texttt{"allServers"} & Omit, or \texttt{"none"} & with Sampling \\
HTTP+SSE transport & Streamable HTTP & sooner \\
\bottomrule
\end{tabularx}
\caption{Deprecated features remain functional. New implementations should not
adopt them. Twelve-month minimum window. See Chapter~\ref{ch:architecture}.}
\label{tab:deprecated-reference}
\end{table}

\section{Version history}

\begin{table}[htbp]
\centering
\small
\begin{tabularx}{\textwidth}{@{}llX@{}}
\toprule
\thead{Version} & \thead{Era} & \thead{Headline} \\
\midrule
\texttt{2024-11-05} & legacy & First release. HTTP+SSE transport \\
\texttt{2025-03-26} & legacy & Streamable HTTP; HTTP+SSE deprecated \\
\texttt{2025-06-18} & legacy & Elicitation, structured output, \texttt{MCP-Protocol-Version} \\
\texttt{2025-11-25} & legacy & URL-mode elicitation, task experiments \\
\texttt{2026-07-28} & \textbf{modern} & Stateless core, MRTR, extensions, lifecycle policy \\
\bottomrule
\end{tabularx}
\caption{``Era'' is the distinction that matters for interoperability: legacy
versions open with \texttt{initialize}. See Chapter~\ref{ch:architecture}.}
\label{tab:version-history}
\end{table}
